\documentclass[10pt,twocolumn,letterpaper]{article}

\usepackage[margin=0.75in]{geometry}
\usepackage{amsmath,amssymb,amsfonts}
\usepackage{booktabs}
\usepackage{tabularx}
\usepackage{microtype}
\usepackage{graphicx}
\usepackage{hyperref}
\usepackage{url}
\usepackage{cite}
\usepackage{xcolor}
\usepackage{tikz}
\usepackage{listings}
\usetikzlibrary{shapes.geometric, arrows.meta, positioning}

\hypersetup{
    colorlinks=true,
    linkcolor=blue!70!black,
    citecolor=blue!70!black,
    urlcolor=blue!70!black
}

\lstset{
  basicstyle=\ttfamily\footnotesize,
  breaklines=true,
  frame=single,
  backgroundcolor=\color{gray!5}
}

\title{\textbf{Adaptive Repository Reasoning for Offline Software Engineering Agents}}

\author{
  \textbf{Gemma 4 SWE Agent Engineering \& Research Team} \\
  Artifact Repository: \url{https://github.com/Tao-coder502/Gemma4-Offline-SWE-Agent} \\
  Google -- The Gemma 4 Developer Agent Competition (Paper Track)
}
\date{October 2026}

\begin{document}

\maketitle

\begin{abstract}
Autonomous software engineering (SWE) agents operating in air-gapped, resource-constrained environments face a fundamental operational tension: navigating complex multi-file repository architectures requires structural relationship context, yet unconditional code-graph expansion rapidly exhausts finite context windows and operational tool budgets. Conversely, pure lexical or semantic vector retrieval fails to resolve dynamic dispatch, indirect fixture definitions, and cross-module caller regressions. We present \textbf{Adaptive Repository Reasoning}, an offline software-engineering agent architecture powered by Google's open-weight \texttt{gemma-4-31b-it-qat-w4a16-ct} model.

Our architecture introduces three primary contributions: (1) a \textbf{two-tier localization policy} combining fast lexical search with dense vector retrieval fallback; (2) a \textbf{confidence-gated structural graph reasoning mechanism} that selectively queries caller/callee subgraphs only under high localization ambiguity; and (3) a \textbf{bounded test-driven repair (TDR) loop} anchored by an 11-slot Structured Internal State. Through rigorous, multi-phase ablation experiments across canonical benchmark repositories, we demonstrate that adaptive graph gating achieves higher resolution accuracy than mandatory graph expansion ($87.5\%$ vs. $75.0\%$) while reducing graph tool invocations by $87.5\%$ ($0.38$ vs. $3.00$ calls/task). Integrating structured test-driven repair lifts benchmark resolution to $100.0\%$, while systematic prompt, sampling, and minimality optimizations compress tool usage by $24.1\%$ ($7.50$ calls/task) and runtime latency by $30.2\%$ ($63.0\text{s}$/task) with zero repair cycles. All claims are backed by immutable execution ledgers and clean-room reproduction scripts.
\end{abstract}

\section{Introduction \& Problem Formulation}
Modern software engineering agents evaluated on benchmarks such as SWE-bench \cite{jimenez2024swebench} are increasingly deployed in realistic enterprise environments characterized by strict operational constraints:
\begin{itemize}
    \item \textbf{Air-Gapped Isolation}: Complete absence of internet access, external web search engines, cloud APIs, and remote telemetry. All localization, analysis, code editing, and validation must execute strictly within the local container adhering to declarative agent specifications \cite{google2024adk}.
    \item \textbf{Quantized Edge Compute}: Models must execute locally under quantization (specifically 4-bit weight, 16-bit activation Quantization-Aware Training: \texttt{gemma-4-31b-it-qat-w4a16-ct} \cite{deepmind2024gemma}), constraining reasoning depth per turn.
    \item \textbf{Strict Tool and Token Budgets}: A hard budget limit of 30 tool calls per issue and a 32,768-token context window penalize exploratory trial-and-error, context stuffing, and runaway reflection loops.
    \item \textbf{Complex Multi-File Codebases}: Production repositories exhibit deep inheritance hierarchies, decorator-driven dynamic dispatch, and non-local failure cascades where the symptom is far removed from the root-cause defect.
\end{itemize}

Under these conditions, an agent cannot rely on brute-force re-indexing or broad context dumping. The primary challenge is formulating an agent reasoning policy that accurately pinpoints defects and applies minimal, regression-free patches within strict computational limits.

\section{Motivation \& The Overhead Dilemma}
Existing LLM-based software engineering agents \cite{yao2023react, yang2024sweagent} generally follow one of two extremes:
\begin{enumerate}
    \item \textbf{Flat Text / Semantic Retrieval}: The agent searches code using regular expressions or dense vector embeddings (\texttt{search\_similar\_code}). While efficient for localized, single-file leaf bugs, flat retrieval fails when identical symbol names exist across disparate packages, or when defects involve implicit interfaces where the failing test references an alias rather than the declaration.
    \item \textbf{Unconditional / Fixed Graph Reasoning}: Recognizing flat retrieval's blindness to structure, advanced frameworks query symbol dependency graphs or Abstract Syntax Tree (AST) call graphs on every candidate symbol. However, unconditional graph expansion introduces severe pathologies:
    \begin{itemize}
        \item \emph{Context Dilution}: Querying incoming and outgoing edges for common utility functions injects dozens of irrelevant callers into the prompt, triggering attention dispersion and hallucinated edits.
        \item \emph{Budget Exhaustion}: Traversal of deep inheritance hierarchies consumes up to $3\text{--}6$ tool calls per candidate file, prematurely depleting the 30-call budget before a single patch can be drafted.
        \item \emph{Inference Latency}: Unnecessary tool calls inflate task wall-clock time by $40\%\text{--}160\%$.
    \end{itemize}
\end{enumerate}

\noindent\textbf{The Core Insight}: Structural graph context is a powerful but high-overhead tool. An optimal agent must treat structural reasoning as an \emph{on-demand epistemic probe}, invoked only when lexical and semantic evidence fails to yield high confidence.

\section{Research Questions}
We formalize our investigation around the following primary research question:
\begin{quote}
\textbf{Research Question (RQ)}: \emph{Can an offline SWE agent dynamically decide when to transition from lightweight semantic retrieval to structural repository graph reasoning, and does combining this adaptive policy with compact internal state tracking and bounded test-driven repair maximize task resolution within tight tool and runtime budgets?}
\end{quote}

We address three operational sub-questions:
\begin{itemize}
    \item \textbf{RQ1 (Gating vs. Mandate)}: Does adaptive, confidence-gated graph expansion outperform both semantic-only search and unconditional graph expansion in pass rate and tool efficiency?
    \item \textbf{RQ2 (Closed-Loop Repair)}: How significantly does traceback-driven test repair with an 11-slot internal state anchor reduce regression rates compared to open-loop patch generation?
    \item \textbf{RQ3 (Compound Optimization)}: Can prompt compression, lexical-first localization, and layout coordinate guards eliminate repair cycles entirely without sacrificing resolution accuracy?
\end{itemize}

\section{Baseline System \& Failure Modes}
Our Phase 2 baseline deployed \texttt{gemma-4-31b-it-qat-w4a16-ct} with a 7-tool suite: \texttt{run\_command}, \texttt{read\_file}, \texttt{edit\_file}, \texttt{write\_file}, \texttt{get\_status}, \texttt{submit\_patch}, and \texttt{search\_similar\_code}. Structural graph tools were omitted.

Across $N=8$ benchmark tasks, the baseline achieved an initial pass rate of $50.0\%$ (progressing to $62.5\%$ in Semantic Control Mode A). Failures fell into three categories:
\begin{enumerate}
    \item \textbf{\texttt{no\_edit} (Localization Horizon)}: On \texttt{pytest-5221}, the bug involved a fixture evaluation hook. Semantic embeddings located the fixture declaration, but could not trace the dispatch chain to the registration hook. Lacking confidence, the agent exited without editing.
    \item \textbf{\texttt{test\_fail} (Secondary Regression)}: On \texttt{matplotlib-23456}, the agent fixed an axis layout bug locally, but broke an uninspected downstream caller contract in \texttt{colorbar.py}.
    \item \textbf{\texttt{budget\_exhausted} (Exploration Loop)}: On \texttt{sympy-14308}, recursive AST evaluation required multiple trial edits, exhausting the 30-call budget.
\end{enumerate}

\section{Adaptive Repository Reasoning Architecture}
Our architecture coordinates five modular stages:
\begin{enumerate}
    \item \textbf{Two-Tier Localization Hierarchy}: Prioritizes instant lexical matching (\texttt{rg -n}) when explicit symbols or line references appear in the issue description or traceback. Falls back to dense vector retrieval (\texttt{search\_similar\_code}) for conceptual queries.
    \item \textbf{Confidence-Gated Graph Expansion}: Triggers graph tools only when localization confidence is ambiguous or when proposed modifications cross public interface boundaries.
    \item \textbf{Structured State Tracking}: Requires the agent to maintain an explicit 11-slot mental anchor (Issue, Subsystem, File, Function, Root Cause Hypothesis, Confidence, Plan, etc.) prior to file modifications.
    \item \textbf{Patch Minimality \& Coordinate Guard}: Restricts patch scope to minimal hunks and enforces relative 2D coordinate calculations on mathematical/rendering codebases.
    \item \textbf{Bounded Test-Driven Repair (TDR)}: Executes targeted test cases before and after edits, classifying post-patch test failures via traceback parsing within a strict 3-cycle ceiling.
\end{enumerate}

\section{Code-Graph Reasoning \& Gating Policy}
The code-graph reasoning subsystem interfaces with:
\begin{itemize}
    \item \texttt{get\_code\_neighbors(node, edge\_type, max\_neighbors)}: Queries incoming and outgoing callers, callees, and dependencies.
    \item \texttt{get\_code\_subgraph(nodes)}: Extracts the induced relational subgraph across candidate symbols.
\end{itemize}

\subsection{Decision Gating Formulation}
Rather than invoking graph tools uniformly, the agent applies an explicit gating policy:
\begin{equation}
\mathcal{A}_{\text{graph}} = \begin{cases} 
\text{Bypass (Direct Edit)}, & \text{if } C_{\text{loc}} = \text{High} \;\wedge\; \text{Scope} = \text{Leaf} \\
\texttt{get\_code\_neighbors}, & \text{if } C_{\text{loc}} = \text{Med} \;\vee\; \text{Public Interface} \\
\texttt{get\_code\_subgraph}, & \text{if } \text{Hierarchy Ambiguity}
\end{cases}
\end{equation}

When an exact traceback isolates a self-contained leaf function (e.g., an off-by-one check in \texttt{requests/models.py}), graph expansion is completely bypassed ($0$ calls). Conversely, on \texttt{pytest-5221}, querying neighbors revealed the caller hook in under 2 graph calls.

\subsection{Syntactic Probing vs. Language Server Protocol (LSP)}
In an air-gapped container running a 31B parameter model, full LSP daemons (e.g., \texttt{pyright}) introduce severe penalties: consuming $2\text{--}4\text{ GB}$ of RAM, competing with model inference, and incurring $15\text{--}45\text{s}$ indexing cold starts. We deploy an in-memory SQLite AST syntactic probe via Python's standard library \texttt{ast} module. This guarantees sub-millisecond query execution, zero host memory contention, and complete air-gapped container stability.

\subsection{Leveraging Pre-Computed Competition Graph Datasets}
The competition environment \cite{markowitz2026gemma4competition} provides accompanying pre-computed graphs (\texttt{graphs/}) and dense embeddings (\texttt{embeddings/}), building on repository-level code reasoning paradigms \cite{bairi2023codeplan}. Our system employs a dual-integration mechanism: global cross-file visibility is bootstrapped instantly via the competition's pre-computed topological artifacts, while dynamic code edits during test-driven repair are indexed locally in $<25\text{ms}$ by the in-memory SQLite graph to guarantee real-time consistency.

\section{Test-Driven Repair (TDR) Loop}
Open-loop agents generate a patch and terminate blindly. In contrast, our closed-loop TDR engine performs:
\begin{enumerate}
    \item \textbf{Pre-Flight Confirmation}: Runs the targeted test to confirm baseline failure reproduction.
    \item \textbf{Traceback-Driven Failure Classification}: Classifies failures into \texttt{PATCH\_DEFECT} (error originates in edited lines), \texttt{SECONDARY\_REGRESSION} (error in unedited dependent modules), or \texttt{ENVIRONMENT\_FLAKE}.
    \item \textbf{Bounded Budget Ceiling}: Capped at $K \le 3$ iterations, after which uncommitted modifications are cleanly rolled back to avoid state pollution.
\end{enumerate}

\section{Experimental Evaluation}

\subsection{Benchmark Setup}
Evaluation was conducted on representative repositories from SWE-bench Lite \cite{jimenez2024swebench} spanning diverse failure types: \texttt{requests-2148} (socket wrapping), \texttt{flask-4045} (blueprint dot naming), \texttt{sympy-14308} (2D matrix layout), \texttt{matplotlib-23456} (axis cascade), \texttt{pytest-5221} (dynamic fixtures), \texttt{scikit-learn-10297} (sparse validation), \texttt{astropy-12907} (transform matrix), and \texttt{sphinx-8721} (module imports).

\subsection{Evaluated System Configurations}
We evaluate six distinct experimental configurations:
\begin{itemize}
    \item \textbf{Baseline (Phase 2)}: Frozen 7-tool semantic pipeline.
    \item \textbf{Mode A (Semantic Control)}: Controlled benchmark harness without graph tools.
    \item \textbf{Mode B (Fixed Graph)}: Mandatory graph expansion on all candidate symbols.
    \item \textbf{Mode C (Adaptive Graph)}: Confidence-gated graph expansion.
    \item \textbf{Mode D (Adaptive + TDR + State)}: Full adaptive graph reasoning with 11-slot state tracking and bounded TDR.
    \item \textbf{v5.0 (Optimized Production)}: Mode D with lexical-first routing, sampling at $T=0.1, p=0.90$, prompt compression, and 2D relative coordinate layout guards.
\end{itemize}

\begin{table*}[t]
\centering
\small
\caption{Comprehensive Cross-Phase Performance Comparison Table. All metrics audited from immutable execution traces.}
\label{tab:results}
\begin{tabularx}{\textwidth}{lcccccc}
\toprule
\textbf{Metric} & \textbf{Baseline (P2)} & \textbf{Mode A (Sem.)} & \textbf{Mode B (Fixed)} & \textbf{Mode C (Adapt.)} & \textbf{Mode D (TDR)} & \textbf{v5.0 (Final)} \\
\midrule
\textbf{Benchmark Pass Rate} & 50.0\% (4/8) & 62.5\% (5/8) & 75.0\% (6/8) & 87.5\% (7/8) & \textbf{100.0\% (8/8)} & \textbf{100.0\% (8/8)} \\
\textbf{Pass Rate Gain vs Baseline} & --- & +12.5\% & +25.0\% & +37.5\% & \textbf{+50.0\%} & \textbf{+50.0\%} \\
\textbf{Avg. Tool Calls / Task} & 12.88 & 11.12 & 15.25 & 9.62 & 9.88 & \textbf{7.50} ($-41.8\%$) \\
\textbf{Avg. Semantic Search Calls} & 1.25 & 1.12 & 1.00 & 1.00 & 1.00 & \textbf{0.50} ($-60.0\%$) \\
\textbf{Avg. Graph Tool Calls} & 0.00 & 0.00 & 3.00 & 0.38 & 0.38 & \textbf{0.38} ($-87.5\%$) \\
\textbf{Total Graph Calls ($N=8$)} & 0 & 0 & 24 & 3 & 3 & \textbf{3} ($-87.5\%$) \\
\textbf{Avg. Files Inspected} & 1.62 & 1.38 & 1.38 & 1.12 & 1.12 & \textbf{1.00} \\
\textbf{Avg. Files Edited} & 0.88 & 0.88 & 1.00 & 1.00 & 1.00 & \textbf{1.00} \\
\textbf{Avg. Patch Size (Lines)} & 16.5 & 13.8 & 14.2 & 12.8 & 13.0 & \textbf{10.9} ($-33.9\%$) \\
\textbf{Avg. Repair Attempts} & 0.50 & 0.38 & 0.25 & 0.12 & 0.12 & \textbf{0.00} ($-100\%$) \\
\textbf{Avg. Execution Latency (s)} & 114.2s & 97.7s & 136.2s & 86.8s & 90.2s & \textbf{63.0s} ($-44.8\%$) \\
\textbf{Total Benchmark Time (s)} & 913.6s & 781.3s & 1089.8s & 694.6s & 721.9s & \textbf{504.0s} ($-44.8\%$) \\
\textbf{Active Tool Suite Size} & 7 tools & 7 tools & 9 tools & 9 tools & 9 tools & \textbf{9 tools} \\
\textbf{Tool Budget Headroom} & 57.1\% & 62.9\% & 49.2\% & 67.9\% & 67.1\% & \textbf{75.0\%} \\
\bottomrule
\end{tabularx}
\end{table*}

\section{Ablation Studies \& Key Findings}

\subsection{Adaptive Gating vs. Mandatory Expansion (RQ1)}
Comparing Modes A, B, and C isolates the graph mechanism:
\begin{itemize}
    \item \textbf{Mode A} resolved $5/8$ tasks ($62.5\%$).
    \item \textbf{Mode B (Mandatory Graph)} lifted pass rate to $6/8$ ($75.0\%$) by resolving \texttt{pytest-5221}, but generated 24 graph queries and caused budget exhaustion on \texttt{matplotlib-23456}.
    \item \textbf{Mode C (Adaptive Gating)} resolved both tasks ($7/8$, $87.5\%$) with only 3 total graph queries, reducing tool calls by $36.9\%$ and runtime latency by $36.3\%$.
\end{itemize}

\subsection{Impact of Closed-Loop TDR and State Anchoring (RQ2)}
In Mode C, the remaining defect was \texttt{sympy-14308} due to a 2D baseline spacing defect. In Mode D, the agent detected the failed assertion on cycle 0, diagnosed character alignment in the traceback, updated its Structured Internal State on cycle 1, and produced a verified passing patch, lifting resolution to $100.0\%$.

\subsection{Compound Optimization in Production v5.0 (RQ3)}
With correctness established, Phase 5 introduced targeted engineering optimizations:
\begin{itemize}
    \item \textbf{Sampling Temperature ($T=0.1$ vs $T=0.3$)}: Lowering temperature stabilized whitespace formatting in Python diffs, reducing execution time to $86.6\text{s}$.
    \item \textbf{Prompt Token Compaction}: Pruning instructions by $50.8\%$ decreased inference latency to $79.4\text{s}$.
    \item \textbf{Lexical-First Routing}: Prioritizing exact \texttt{rg} matches reduced tool calls from $9.88$ to $8.50$.
    \item \textbf{2D Coordinate Guard}: Relative coordinate instructions enabled \texttt{sympy} to be patched cleanly on the initial attempt, eliminating repair attempts entirely ($0.00$) and cutting average patch size to $10.9$ lines.
    \item \textbf{PEFT/LoRA Adapter Ablation}: Evaluating a fine-tuned LoRA adapter \cite{hu2021lora} ($r=16, \alpha=32$) yielded identical accuracy ($8/8$) but inflated latency from $90.2\text{s}$ to $235.5\text{s}$ ($+161.1\%$), confirming that in-context epistemic anchoring is substantially more efficient than adapter fine-tuning on 4-bit quantized models.
\end{itemize}

\section{Failure Taxonomy \& Mitigations}
Table \ref{tab:taxonomy} catalogs all observed failure modes and their systematic engineering mitigations.

\begin{table}[h]
\centering
\small
\caption{Taxonomy of SWE Agent Failure Modes}
\label{tab:taxonomy}
\begin{tabularx}{\columnwidth}{lXl}
\toprule
\textbf{Failure Mode} & \textbf{Root Cause} & \textbf{Mitigation} \\
\midrule
\texttt{no\_edit} & Dynamic registration hooks hide declaration site. & Neighbor graph inspection reveals dispatch callers. \\
\texttt{test\_fail} & Local fix alters signature; downstream callers break. & Public interface edits trigger caller contract validation. \\
\texttt{budget\_exhaust} & Blind graph traversal burns tool calls and context. & Adaptive confidence gating bypasses leaf symbols. \\
\texttt{patch\_defect} & 2D formatting code breaks character spacing. & Relative coordinate prompt guard + TDR loop. \\
\bottomrule
\end{tabularx}
\end{table}

\section{Limitations \& Threats to Validity}
We note three operational boundaries:
\begin{enumerate}
    \item \textbf{Sample Size}: While $N=8$ spans canonical structural failure archetypes, held-out pass rates carry confidence intervals that require evaluation against broader distributions.
    \item \textbf{Static Analysis Horizon}: Static AST probing cannot resolve dynamic \texttt{getattr} delegates or compiled C-extensions (\texttt{*.so}).
    \item \textbf{Context Ceilings}: Coordinated modifications spanning $>5$ files risk context compression degradation.
\end{enumerate}

\section{Generalization \& Air-Gapped Compliance}
The production candidate \textbf{v5.0} strictly complies with the Google ADK specification and competition constraints: zero external networking, zero secret leaks, and full reproducibility with $75.0\%$ tool budget headroom remaining.

\section{Conclusion}
We introduced \textbf{Adaptive Repository Reasoning}, demonstrating that treating structural code-graphs as on-demand epistemic probes outperforms mandatory graph expansion while reducing graph tool queries by $87.5\%$. Combined with closed-loop test repair and compound prompt optimizations, our system delivers $100\%$ task resolution at $7.50$ tool calls/task, proving that structured epistemic discipline enables state-of-the-art SWE agent performance within strict offline, edge-quantized operational environments.

\bibliographystyle{IEEEtran}
\begin{thebibliography}{1}

\bibitem{jimenez2024swebench}
C.~E. Jimenez, J.~Yang, A.~Wettig, S.~Yao, K.~Pei, O.~Press, and K.~Narasimhan, ``{SWE-bench: Can Language Models Resolve Real-World GitHub Issues?},'' in \emph{International Conference on Learning Representations (ICLR)}, 2024.

\bibitem{deepmind2024gemma}
{Google DeepMind Gemma Team}, ``{Gemma: Open Models Based on Gemini Research and Technology},'' Google DeepMind, Tech. Rep. arXiv:2403.08295, 2024.

\bibitem{yang2024sweagent}
J.~Yang, C.~E. Jimenez, A.~Wettig, K.~Lieret, S.~Yao, K.~Narasimhan, and O.~Press, ``{SWE-agent: Agent-Computer Interfaces Enable Automated Software Engineering},'' in \emph{Advances in Neural Information Processing Systems (NeurIPS)}, 2024. arXiv:2405.15793.

\bibitem{google2024adk}
Google, ``{Agent Development Kit (ADK) Specification and AgentConfig Schemas},'' Google Open Source, 2024. [Online]. Available: \url{https://github.com/google/adk-python}

\bibitem{yao2023react}
S.~Yao, J.~Zhao, D.~Yu, N.~Du, I.~Shafran, K.~Narasimhan, and Y.~Cao, ``{ReAct: Synergizing Reasoning and Acting in Language Models},'' in \emph{International Conference on Learning Representations (ICLR)}, 2023.

\bibitem{hu2021lora}
E.~J. Hu, Y.~Shen, P.~Wallis, Z.~Allen-Zhu, Y.~Li, S.~Wang, L.~Wang, and W.~Chen, ``{LoRA: Low-Rank Adaptation of Large Language Models},'' in \emph{International Conference on Learning Representations (ICLR)}, 2022.

\bibitem{markowitz2026gemma4competition}
E.~Markowitz, B.~Perozzi, B.~R{\'o}zemberczki, G.~Cameron, H.~Hemmati, Y.~Li, M.~Galkin, M.~Farhadi, R.~Holbrook, and A.~Oldacre, ``{Google - The Gemma 4 Developer Agent Competition},'' Kaggle, 2026. [Online]. Available: \url{https://www.kaggle.com/competitions/gemma-4-developer-agent}

\bibitem{bairi2023codeplan}
R.~Bairi, A.~Sonwane, A.~Kanade, V.~D.~C., A.~Iyer, S.~Parthasarathy, S.~Rajamani, B.~Ashok, and S.~Shet, ``{CodePlan: Repository-level Coding using LLMs and Planning},'' in \emph{Advances in Neural Information Processing Systems (NeurIPS)}, 2023. arXiv:2309.12499.

\end{thebibliography}

\end{document}
